\uuid{mH6z}
\titre{Convolution 2D complète à la main avec padding same}
\niveau{L3}
\module{Informatique / Signal}
\chapitre{Réseaux de neurones}
\sousChapitre{Convolution 2D}
\theme{Sobel X, padding same, noyau antisymétrique}
\auteur{}
\datecreate{2026-03-23}
\organisation{}
\difficulte{3}

\contenu{
	\texte{
		On considère l'image $I$ et le noyau $K$ suivants :
		\[
		I = \begin{pmatrix}
			1 & 0 & 1 & 2 \\
			3 & 1 & 0 & 1 \\
			2 & 2 & 1 & 0 \\
			0 & 1 & 3 & 1
		\end{pmatrix},
		\qquad
		K = \begin{pmatrix}
			1 & 0 & -1 \\
			2 & 0 & -2 \\
			1 & 0 & -1
		\end{pmatrix}.
		\]
		On applique ce filtre en mode \textbf{same} (padding $p = 1$) avec stride $= 1$.
	}
	\begin{enumerate}
		\item
		\question{Écrire l'image paddée $I_p$ ($6 \times 6$ avec des zéros sur les bords).}
		\reponse{
			\[
			I_p = \begin{pmatrix}
				0 & 0 & 0 & 0 & 0 & 0 \\
				0 & 1 & 0 & 1 & 2 & 0 \\
				0 & 3 & 1 & 0 & 1 & 0 \\
				0 & 2 & 2 & 1 & 0 & 0 \\
				0 & 0 & 1 & 3 & 1 & 0 \\
				0 & 0 & 0 & 0 & 0 & 0
			\end{pmatrix}.
			\]
		}
		
		\item
		\question{Rappeler le noyau renversé $K' = \mathrm{flip}(K)$. Que remarque-t-on ?}
		\reponse{
			Le retournement (rotation de $180^\circ$) de $K$ donne :
			\[
			K' = \begin{pmatrix}
				-1 & 0 & 1 \\
				-2 & 0 & 2 \\
				-1 & 0 & 1
			\end{pmatrix}.
			\]
			On remarque que $K' = -K$ : le noyau est \textbf{antisymétrique par rotation de
				$180^\circ$}. Cela implique que la convolution $I \star K$ et la corrélation croisée
			$I \star\star K$ diffèrent d'un signe.
		}
		
		\item
		\question{Calculer les 4 coefficients de la ligne du milieu de la sortie,
			aux positions $(1,0)$, $(1,1)$, $(1,2)$, $(1,3)$ (indexation ligne/colonne à partir de 0).}
		\reponse{
			On travaille sur $I_p$ et on extrait les fenêtres $3 \times 3$ centrées sur les
			pixels $(1,0)$ à $(1,3)$ de $I$ (ce qui correspond aux lignes d'indices 1 à 3 de $I_p$
			grâce au décalage de $p=1$). Le produit scalaire se fait avec $K'$.
			
			\textbf{Position $(1,0)$} — fenêtre $I_p[1:4,\; 0:3]$ :
			\[
			\begin{pmatrix}0&1&0\\0&3&1\\0&2&2\end{pmatrix} \odot
			\begin{pmatrix}-1&0&1\\-2&0&2\\-1&0&1\end{pmatrix}
			= 0+0+0 + 0+0+2 + 0+0+2 = 4.
			\]
			
			\textbf{Position $(1,1)$} — fenêtre $I_p[1:4,\; 1:4]$ :
			\[
			\begin{pmatrix}1&0&1\\3&1&0\\2&2&1\end{pmatrix} \odot K'
			= (-1)+0+1 + (-6)+0+0 + (-2)+0+1 = -7.
			\]
			
			\textbf{Position $(1,2)$} — fenêtre $I_p[1:4,\; 2:5]$ :
			\[
			\begin{pmatrix}0&1&2\\1&0&1\\2&1&0\end{pmatrix} \odot K'
			= 0+0+2 + (-2)+0+2 + (-2)+0+0 = 0.
			\]
			
			\textbf{Position $(1,3)$} — fenêtre $I_p[1:4,\; 3:6]$ :
			\[
			\begin{pmatrix}1&2&0\\0&1&0\\1&0&0\end{pmatrix} \odot K'
			= (-1)+0+0 + 0+0+0 + (-1)+0+0 = -2.
			\]
			
			Ligne du milieu de la sortie : $\begin{pmatrix}4 & -7 & 0 & -2\end{pmatrix}$.
		}
		
		\item
		\question{Interpréter le résultat : ce filtre est-il symétrique ? Que détecte-t-il ?}
		\reponse{
			$K$ n'est pas symétrique : il vaut $+1$ à gauche et $-1$ à droite (colonnes).
			C'est le filtre de \textbf{Sobel X} : il calcule approximativement la dérivée
			partielle horizontale $\partial I / \partial x$.
			Une valeur positive indique une transition sombre $\to$ clair de gauche à droite~;
			une valeur négative indique clair $\to$ sombre. Il détecte donc les
			\textbf{contours verticaux} (variations horizontales d'intensité).
		}
	\end{enumerate}
}